\subsection{Objeto Virtual de aprendizaje (OVA)}
El Ministerio de Educación Nacional – MEN, define un objeto de aprendizaje como: ``un conjunto de recursos digitales, autocontenible y reutilizable, con un propósito educativo y constituido por al menos tres componentes internos: contenidos, actividades de aprendizaje y elementos de contextualización. El Objeto de Aprendizaje debe tener una estructura de información externa (metadatos) que facilite su almacenamiento, identificación y recuperación" \cite{men0}.

\ \\  
\textbf{Elementos estructurales de un Objeto de Aprendizaje} \cite{men0}.
\begin{itemize}
\item \textbf{Objetivos:} Expresan de manera explícita lo que el estudiante va a aprender.
\item \textbf{Contenidos:} Se refiere a los conceptos y sus múltiples formas de representarlos, pueden ser: definiciones, entrevistas, imágenes, etc.
\item \textbf{Actividades de Aprendizaje:} Tareas que el estudiante debe realizar con base en el tema estudiado, con el fin de hacer significativo el aprendizaje, desarrollar habilidades y alcanzar los objetivos de aprendizaje.
\item \textbf{Elementos de contextualización:} Hace referencia a los datos que describen el objeto, como: título, idioma, la versión, la información relacionada con los derechos de autor.
\end{itemize}